\chapter{自演化数字孪生}
\label{ch:evolve}

单次烹饪开始时，系统只有该食材类别的先验。结束时，它应当留下对这一类别、这一厚度区间仍然有效的参数后验。演化发生在参数与算子修正项上，不发生在能量方程的结构上。

\section{离线算子与在线更新的分工}

全场方程 \eqref{eq:energy}--\eqref{eq:structure} 的求解用于离线生成轨迹族。在线预测使用
\begin{equation}
\label{eq:op}
\vect{z}_{k+1}=\mathcal{G}_{\phi}(\vect{z}_k,\vect{u}_k,\vect{\theta})+\vect{r}_{\psi}(\vect{z}_k,\vect{u}_k),
\end{equation}
$\mathcal{G}_{\phi}$ 是在仿真轨迹上训练的神经算子。Fourier Neural Operator 与 DeepONet 使用同一训练集，主模型取蒸汽干预位移误差更低者，另一套保留为对照。$\vect{\theta}$ 是物理参数，$\vect{r}_{\psi}$ 是残差修正。算子输入包含执行器轨迹，因此干预是替换输入后的同一次前向计算。

残差约束为
\begin{equation}
\label{eq:bnp}
\|\vect{r}_{\psi}\|\le r_{\max}(\vect{z},\vect{u}),
\end{equation}
$r_{\max}$ 由离线验证集上的误差分位数确定。超出该界的网络输出被截断，并回退到两节点物理模型。这是“物理先验不会被一次异常烹饪写坏”的具体机制。

\section{贝叶斯更新}

类别先验 $p(\vect{\theta}\mid\mathrm{class})$ 由离线仿真加公开的热物性范围构成。窗 $w$ 内的似然来自 \eqref{eq:obs}。更新
\begin{equation}
\label{eq:bayes}
p(\vect{\theta}\mid\vect{y}_{0:w},\mathrm{class})
\propto p(\vect{y}_{0:w}\mid\vect{\theta})\,p(\vect{\theta}\mid\mathrm{class}).
\end{equation}
在线用集合或变分高斯近似这个后验。烹饪结束时，把充分统计量（均值、协方差、食材类别、厚度区间、探针是否有效）写回知识库。下一次同类别的先验是这次后验与类别超先验的精度加权，而不是用一条轨迹覆盖整个类别。

\begin{equation}
\label{eq:hier}
\vect{\Lambda}_{\mathrm{next}}=\vect{\Lambda}_{\mathrm{hyper}}+\sum_{j\in\mathrm{class}}\vect{\Lambda}_{j},
\qquad
\vect{\Lambda}_{\mathrm{next}}\vect{\mu}_{\mathrm{next}}
=\vect{\Lambda}_{\mathrm{hyper}}\vect{\mu}_{\mathrm{hyper}}+\sum_{j}\vect{\Lambda}_{j}\vect{\mu}_{j}.
\end{equation}
厨房之间的差异不会在本项目里上传。知识库是本地的。

\section{一次烹饪内部的演化}

结壳使 $R_{\mathrm{int}}$ 在几十分钟内漂移，漂移比类别差异更快。过程模型
\begin{equation}
R_{\mathrm{int},k+1}=R_{\mathrm{int},k}+\kappa_c\langle M_c-\varepsilon_{l,s}\rangle_+\Delta t+\xi_k
\end{equation}
与 \eqref{eq:structure} 一致。粒子滤波同时跟踪 $R_{\mathrm{int}}$ 与 $T_2$。探针 P4 在漂移后可以再触发一次，触发间隔不短于 $t_{\mathrm{gap}}$，避免探针本身成为主要加热方式。

\section{算子训练数据如何覆盖干预}

训练集如果只包含 PID 控温轨迹，神经算子在蒸汽阶跃上会外推失败，因果干预就不可信。数据生成协议：

\begin{enumerate}
\item 食材：瘦肉糜板块、土豆块、面团，三种孔隙与反应不同的代表；
\item 几何：厚度 3 档、初始温度 2 档、初始含水 2 档；
\item 执行器：加热、蒸汽、风速、排汽的拉丁超立方，外加探针库 P1--P4 的插入；
\item 每条轨迹保存真值 $\vect{z}(t)$、可安装传感器的带噪 $\vect{y}(t)$、执行器 $\vect{u}(t)$；
\item 划分按几何厚度分组，禁止同一厚度的相邻轨迹同时出现在训练和干预测试中。
\end{enumerate}

教师轨迹由高保真一维或二维有限元生成，腔体边界来自少量 CFD 或关联式。高保真数量为数百条。降阶模型可把轨迹扩充到 $10^4$ 量级，扩充条数是数值样本，不是独立实验。每条轨迹的方程、网格和扩充规则写入数据索引。

\section{失效与回退}

\begin{itemize}
\item 物理残差在滚动窗口上超过 $r_{\max}$：丢弃 $\vect{r}_{\psi}$，只用 $\mathcal{G}$ 的物理降阶部分；
\item 后验协方差比先验更大：本次不写回知识库；
\item 类别识别置信度低：使用所有类别中最保守的 $\alpha_m$ 分位数，并提高探针倾向。
\end{itemize}

更新规则因此是三条硬限制。在线只动预先声明的传递参数和残差项。残差超出离线误差分位界时，残差项丢弃，预测回到物理降阶模型。单次过程只有在后验精度高于先验时，才按精度加权写入本地类别分布。

算子损失含三项：轨迹误差、干预位移误差、残差越界惩罚。层次更新的精度公式即 \eqref{eq:hier}。回退是否被触发，由第 \ref{ch:sim} 章的残差日志计数。
